 \begin{propertie} \textbf{Algebra de Integral de l\'inea.}\label{alge} Sea $C^{+}$ una curva simple, orientada y suave a trozos, adem\'as,  sean  $\mathbf{F}$ y $\mathbf{G}$ campos vectoriales continuos con dominios tales que incluyen a la curva $C$, entonces: 
\begin{itemize}
\item[1.]  \[
        \int_{C^{+}}(\mathbf{F} + \mathbf{G})  \cdot d\mathbf{s}=  \int_{C^{+}} \mathbf{F} \cdot d\mathbf{s}  +  \int_{C^{+}} \mathbf{G}  \cdot d\mathbf{s} 
    \]

\item[2.]  \[
        \int_{C^{+}} k \mathbf{F}  \cdot d\mathbf{s}= k \int_{C^{+}} \mathbf{F} \cdot d\mathbf{s}   
    \]

\item[3.]  \[
        \int_{C^{+}} k \mathbf{F}  \cdot d\mathbf{s}= - \int_{C^{-}} \mathbf{F} \cdot d\mathbf{s}   
    \]

\item[3.]  Si $C = \cup_{i=1}^{n} C_{i},$ entonces \[
        \int_{C^{+}}  \mathbf{F}  \cdot d\mathbf{s}=  \sum_{i=1}^{n}  \int_{C_{i}^{+}} \mathbf{F} \cdot d\mathbf{s}   
    \]

\end{itemize}
 \end{propertie}


 \begin{nota}\label{not_int} Sean  $\mathbf{F}=(f_1, f_2)$ y $\mathbf{G}=(g_1, g_2, g_3)$ campos vectoriales, notaremos $$ \int_{C_1}\mathbf{F}\cdot d\mathbf{s} =  \int_{ C_1}f_1\:dx+ f_2 \:dy$$ 
  $$ \int_{C_2}\mathbf{G}\cdot d\mathbf{s} =  \int_{C_2} g_1\:dx+g_2\:dy+g_3\:dz.$$
 \end{nota}
 

\begin{nota} Si $C$  es una curva simple cerrada,  se suele notar a la integral como
    \[
        \oint_{C}\mathbf{F}\cdot d\mathbf{s},  
    \]  y se la llama integral de circulaci\'on de $\mathbf{F}$ sobre $C$.
\end{nota}




\begin{example} Si $C$ es la curva simple frontera de la región delimitada por la recta $y = x$ y la parábola $y = 6x - x^2$, calcular, indicando la orientaci\'on elegida,   la integral de línea $$\displaystyle\oint_{C} xy\,dx + y^2\,dy.$$ 


\textbf{Solución:} Las curvas que delimitan la región se cortan cuando $x = 6x - x^2$, es decir $x(x-5) = 0$, lo que da los puntos $A(0,0)$ y $B(5,5)$. Para $0 \le x \le 5$ se verifica $x \le 6x - x^2$, de modo que la parábola queda por encima de la recta y la región encerrada por $C$ es
\[
D = \{(x,y)\in\mathbb{R}^2 : 0 \le x \le 5,\ x \le y \le 6x - x^2\}.
\]

\begin{center}
\begin{tikzpicture}[scale=0.55]
    % Ejes
    \draw[->] (-0.5, 0) -- (7, 0) node[right] {$x$};
    \draw[->] (0, -0.5) -- (0, 10.5) node[above] {$y$};

    % Región sombreada entre la recta y la parábola en [0,5]
    \fill[blue!12, domain=0:5, samples=80] 
        plot ({\x}, {6*\x - \x*\x}) -- (5,5) -- (0,0) -- cycle;

    % Parábola y = 6x - x^2  (tramo C_2, orientado de B a A)
    \draw[thick, domain=0:6, samples=80] plot ({\x}, {6*\x - \x*\x});
    % Flecha de orientación sobre la parábola (de derecha a izquierda)
    \draw[thick, ->] plot [domain=4:3.5, samples=20] ({\x}, {6*\x - \x*\x});

    % Recta y = x  (tramo C_1, orientado de A a B)
    \draw[thick] (0,0) -- (5,5);
    \draw[thick, ->] (2.4, 2.4) -- (2.7, 2.7);

    % Puntos
    \filldraw (0,0) circle (2pt) node[below left] {$A(0,0)$};
    \filldraw (5,5) circle (2pt) node[right] {$B(5,5)$};

    % Etiquetas
    \node at (2.5, 4.3) {$D$};
    \node[below right] at (3.5, 3.5) {$C_1$};
    \node[above] at (1.8, 8.2) {$C_2$};
    \node[right] at (5.2, 9) {$y = 6x - x^2$};
    \node[right] at (5, 5) {};

    % Marcas en los ejes
    \foreach \x in {1,2,3,4,5,6}
        \draw (\x, 0.12) -- (\x, -0.12) node[below] {\small\x};
    \foreach \y in {2,4,6,8}
        \draw (0.12, \y) -- (-0.12, \y) node[left] {\small\y};
\end{tikzpicture}
\end{center}

Elegimos la \emph{orientación positiva} (antihoraria), en la que $D$ queda a la izquierda al recorrer $C$. Con esta convención, la curva se descompone en dos tramos independientes: el segmento $C_1$ que va de $A(0,0)$ hasta $B(5,5)$ sobre la recta $y=x$, y el arco $C_2$ que vuelve de $B(5,5)$ hasta $A(0,0)$ sobre la parábola. Entonces,el trabajo del campo $\mathbf{F}(x,y) = (xy,\, y^2)$ sobre la curva $C$, resulta
\[
\oint_C xy\, dx + y^2\, dy \;=\; \int_{C_1} \mathbf{F}\cdot d\mathbf{s} \;+\; \int_{C_2} \mathbf{F}\cdot d\mathbf{s}.
\]

\textbf{1. Tramo $C_1$:} Desde $(0,0)$ hasta $(5,5)$.
Parametrizamos el segmento lineal mediante la trayectoria $\boldsymbol{\sigma}_1(t) = (t, t)$ con $t \in [0,5]$, cuyo vector tangente es $\boldsymbol{\sigma}_1'(t) = (1, 1)$. Planteando la integral de trayectoria obtenemos:
\[
\int_{C_1} \mathbf{F}\cdot d\mathbf{s} \;=\; \int_0^5 (t^2,\, t^2)\cdot(1,\,1)\, dt \;=\; \int_0^5 2t^2\, dt \;=\; \frac{250}{3}.
\]

\textbf{2. Tramo $C_2$:} Desde $(5,5)$ hasta $(0,0)$.
Parametrizamos el arco de parábola mediante $\boldsymbol{\sigma}_2(t) = (t,\, 6t - t^2)$ recorrido en sentido decreciente, con $t$ variando de $5$ a $0$. Su vector tangente es $\boldsymbol{\sigma}_2'(t) = (1,\, 6 - 2t)$, por lo que
\begin{align*}
\int_{C_2} \mathbf{F}\cdot d\mathbf{s} &= \int_5^0 \big(t(6t - t^2),\, (6t-t^2)^2\big)\cdot(1,\, 6-2t)\, dt \\
&= -\int_0^5 \Big[(6t^2 - t^3) + (6t - t^2)^2(6 - 2t)\Big]\, dt \\
&= -\int_0^5 \big(222\,t^2 - 145\,t^3 + 30\,t^4 - 2\,t^5\big)\, dt \;=\; -\frac{1625}{12}.
\end{align*}

Sumando ambas contribuciones se obtiene finalmente
\[
\oint_C xy\, dx + y^2\, dy \;=\; \frac{250}{3} - \frac{1625}{12} \;=\; \frac{1000 - 1625}{12} \;=\; -\frac{625}{12}.
\]
\end{example}


\begin{example}\label{ej:trabajo_poligonal}
Hallar el trabajo realizado por el campo vectorial $\mathbf{F}(x,y) = (x^2y,\, x-y)$ al mover una partícula a lo largo de la trayectoria poligonal cerrada $C$ que une los puntos $A(0,0) \to B(2,3) \to C(4,2) \to A(0,0)$.
\end{example}

\begin{figure}[htbp]
    \centering
    % Llamamos al archivo con el código TikZ extraído
    \tikzexternaldisable
\begin{center}
\begin{tikzpicture}[scale=1.1, >=stealth]

  % Grilla de fondo sutil
  \draw[very thin, gray!20] (-0.5,-0.5) grid (4.5,3.5);

  % Ejes Cartesianos
  \draw[->, thick] (-0.5, 0) -- (4.8, 0) node[right] {$x$};
  \draw[->, thick] (0, -0.5) -- (0, 3.8) node[above] {$y$};

  % Marcas en los ejes
  \foreach \x in {1,2,3,4}
    \draw (\x, 0.08) -- (\x, -0.08) node[below, font=\small] {$\x$};
  \foreach \y in {1,2,3}
    \draw (0.08, \y) -- (-0.08, \y) node[left, font=\small] {$\y$};

  % Relleno de la región encerrada por la curva C
  \fill[blue!15, opacity=0.7] (0,0) -- (2,3) -- (4,2) -- cycle;

  % Trayectoria poligonal C con orientación
  % Segmento AB (C1)
  \draw[line width=1.1pt, blue!70!black, postaction={decorate, decoration={
    markings, mark=at position 0.55 with {\arrow{>}}}}] (0,0) -- (2,3) 
    node[midway, above left, font=\small] {$C_1$};

  % Segmento BC (C2) -> Se cambió 'above right' por 'above, sloped' o desplazamiento fino
  \draw[line width=1.1pt, blue!70!black, postaction={decorate, decoration={
    markings, mark=at position 0.55 with {\arrow{>}}}}] (2,3) -- (4,2) 
    node[midway, above right=1pt, font=\small] {$C_2$};

  % Segmento CA (C3)
  \draw[line width=1.1pt, blue!70!black, postaction={decorate, decoration={
    markings, mark=at position 0.55 with {\arrow{>}}}}] (4,2) -- (0,0) 
    node[midway, below, font=\small] {$C_3$};

  % Puntos/Vértices A, B, C
  \fill[blue!80!black] (0,0) circle (2pt) node[below left, font=\small\bfseries] {$A(0,0)$};
  \fill[blue!80!black] (2,3) circle (2pt) node[above, font=\small\bfseries] {$B(2,3)$};
  \fill[blue!80!black] (4,2) circle (2pt) node[right, font=\small\bfseries] {$C(4,2)$};

  % Etiquetas D y C
  \node[color=blue!80!black, font=\large] at (2, 2) {$D$};
  % La etiqueta C de la frontera total se coloca cerca del tramo C1 o arriba
  \node[color=blue!70!black, font=\large] at (0.9, 2.5) {$C$};

\end{tikzpicture}
\end{center}
\tikzexternalenable 
    \label{fig:ilustracion ejemplo}
\end{figure}


\noindent\textbf{Solución:} De acuerdo con la definición de integral sobre una trayectoria a trozos, el trabajo total $W$ se obtiene dividiendo la curva en tres tramos independientes ($C = C_{AB} \cup C_{BC} \cup C_{CA}$) y sumando sus respectivas contribuciones: 
\[ 
W = \int_{C} \mathbf{F} \cdot d\mathbf{s} = \int_{C_{AB}} \mathbf{F} \cdot d\mathbf{s} + \int_{C_{BC}} \mathbf{F} \cdot d\mathbf{s} + \int_{C_{CA}} \mathbf{F} \cdot d\mathbf{s}. 
\] 

\noindent\textbf{1. Tramo $C_{AB}$: Desde $(0,0)$ hasta $(2,3)$.} \\ 
Parametrizamos el segmento lineal mediante la trayectoria $\boldsymbol{\sigma}_1(t) = (2t,\, 3t)$ con $t\in[0,1]$, cuyo vector tangente es $\boldsymbol{\sigma}_1'(t) = (2, 3)$. Planteando la integral de trayectoria obtenemos:
\[ 
W_{AB} = \int_0^1 \bigl( (2t)^2(3t) ,\, 2t-3t \bigr) \cdot (2,3) \,dt = \int_0^1 (24t^3 - 3t)\,dt = \frac{9}{2}. 
\] 

\noindent\textbf{2. Tramo $C_{BC}$: Desde $(2,3)$ hasta $(4,2)$.} \\ 
La forma estándar de parametrizar este segmento es $\boldsymbol{\sigma}_2(t) = (2+2t,\, 3-t)$ con $t\in[0,1]$, lo que implica que $\boldsymbol{\sigma}_2'(t) = (2, -1)$. Sustituyendo en el producto escalar vectorial resulta:
\[ 
W_{BC} = \int_0^1 \bigl( (2+2t)^2 (3-t) ,\, (2+2t)-(3-t) \bigr) \cdot (2,-1) \,dt = \int_0^1 (25 + 37t + 8t^2 - 8t^3)\,dt.
\] 
Evaluando la primitiva mediante la regla de Barrow:
\[
W_{BC} = \left[25t + \frac{37}{2}t^2 + \frac{8}{3}t^3 - 2t^4\right]_0^1 = \frac{265}{6}. 
\]

\noindent\textbf{3. Tramo $C_{CA}$: Desde $(4,2)$ hasta $(0,0)$.} \\ 
Parametrizamos el regreso al origen mediante $\boldsymbol{\sigma}_3(t) = (4-4t,\, 2-2t)$ con $t\in[0,1]$, de donde se obtiene $\boldsymbol{\sigma}_3'(t) = (-4, -2)$. Planteamos la integral correspondiente:
\[ 
W_{CA} = \int_0^1 \bigl( (4-4t)^2 (2-2t) ,\, (4-4t)-(2-2t) \bigr) \cdot (-4,-2) \,dt = \int_0^1 \bigl[-128(1-t)^3 - 4(1-t)\bigr]\,dt.
\] 
Resolviendo la integral por el método de sustitución:
\[
W_{CA} = \left[\frac{128(1-t)^4}{4} + \frac{4(1-t)^2}{2}\right]_0^1 = 0 - (32 + 2) = -34. 
\]

\noindent Sumando el trabajo neto realizado por el campo en cada uno de los tres tramos independientes de la poligonal, obtenemos el resultado final: 
\[ 
W = W_{AB} + W_{BC} + W_{CA} = \frac{9}{2} + \frac{265}{6} - 34 = \frac{27 + 265 - 204}{6} = \frac{88}{6} = \frac{44}{3}. \tag*{$\blacksquare$} 
\] 




 \begin{tarea} Sea $C$ la curva simple dada por la uni\'on de los segmentos que unen los puntos $(0,0)$,  $(0,3)$ y $(0,2)$. Calcular, indicando la orientaci\'on elegida para $C$  $$ \int_{C}  (y^{2}, 2xy +1)  \cdot d\mathbf{s} $$
 \end{tarea}

 \begin{tarea} Mismo ejercicio con la orientaci\'on para $C$ contraria  a la elegida y comparar los resultados obtenidos. 
 \end{tarea}





